\clearpage
\subsection{题目2.3.1：正数点积测试}
两组有效输入均为[1,2,3,4,5,6,7,8]，八项平方和为204。历史全核functional-2记录中，seq=1、ROB=0:39、pdest=103在339296周期接受、339298周期响应；实际完整128位结果为\seqsplit{0x000000000000000000000000000000cc}，与预期相同。源高64位为A5哨兵，不参与点积。下图给出对应输入与实际响应摘录。\par
\reportfigure{S2-TEST-01}{正数点积测试：预期与实际结果均为204}{5.0cm}{两个源的有效元素均为1至8；expected与actual完整128位一致，末尾为CC。}{用户保存的原始PNG；functional-2观察器日志字段摘录，保留事务身份、原始行号与SHA256；不是波形截图。}
\subsection{题目2.3.2：负数点积测试}
有效输入为八个−1与八个1，预期总和−8。functional-0的seq=4、ROB=0:12、pdest=143在420291周期接受、420293周期响应；实际完整结果为\seqsplit{0x000000000000000000000000fffffff8}。低32位保存负数补码，高96位为零；该检查不同于第一阶段按字节相加的vadd补码自检。\par
\reportfigure{S2-TEST-02}{负数点积测试：低32位为−8的补码，高96位清零}{5.0cm}{八个−1与八个1点积；实际低32位为FFFFFFF8，高96位为零。}{用户保存的原始PNG；functional-0观察器日志字段摘录，保留事务身份、原始行号与SHA256；不是波形截图。}
\clearpage
\subsection{题目2.3.3：边界值点积测试}
最大正结果使用八个−128与八个−128，每项乘积16384，总和131072。functional-0的seq=5、ROB=0:43、pdest=66在448417周期接受、448419周期响应；实际完整结果为\seqsplit{0x00000000000000000000000000020000}。该例检验乘法和归约在正端点不被截断。\par
\reportfigure{S2-TEST-03}{边界值点积测试：最大正结果131072}{5.0cm}{八对−128相乘求和；实际低32位为00020000，高96位为零。}{用户保存的原始PNG；functional-0观察器日志字段摘录，保留事务身份、原始行号与SHA256；本图只证明所示最大正结果案例。}
零输入用八个0与八个127，实际完整结果为零；functional-0的seq=0、ROB=0:8、pdest=124在311084周期接受、311086周期响应。最小结果的数学预期为8×(−128)×127=−130048，低32位为\seqsplit{0xfffe0400}；本稿未单独展开该项全核输入与响应配对，因此该值只列为预期边界，不新增通过声明。\par
\noindent\textbf{证据来源与判读：}固定0f历史运行，原始记录根目录为\seqsplit{/home/riscv/riscv-lab/software-import-20261026/clean-02/fullcore-run-01/runs/}；正数在\seqsplit{functional-2/run.log}第7030—7034行，负数、最大正结果、零输入分别在\seqsplit{functional-0/run.log}第7927—7931、8400—8404、5814—5818行。按seq、ROB与pdest配对，比较完整128位响应。上述三图是日志字段摘录，不是GTKWave截图；执行响应不等同于架构提交，完整运行终态与差分结论按原归档核对。\par

\subsection{算术数据通路测试汇总}
除前述具体用例外，历史组合算术RTL测试统计为22项定向测试、10000项随机测试，合计10022项全部通过。该汇总与前述全核用例摘录属于不同验证层级，各自保留原测试范围。\par
\reportfigure{S2-TEST-SUMMARY}{VDot算术数据通路定向与随机测试结果}{4.0cm}{定向测试22项、随机测试10000项，合计10022项全部通过。本结果属于算术数据通路验证；全核协同仿真、流水线反压及取消行为由相应测试另行验证。}{用户提供的原始统计截图，完整保留三行PASS输出。本次仅补图，未重新执行测试；后续S2-09综合图中的算术统计为同一组历史结果，不重复计数。}
